\section{Zadanie 6}
\autor{Mateusz Wojaczek}
Udowodnij, że dla dowolnego skończonego nieprzywiedlnego łańcucha Markowa
istnieją stałe \(c > 0\), \(0 < \delta < 1\) takie, że dla każdego stanu \(z\) i dla każdego \(m \in \mathbb{N}_1\) mamy:
\[
	\mathbb{P}(T_z > m) \leq c \, \delta^m.
\]

\subsection*{Rozwiązanie}
Udowodnij, że dla dowolnego skończonego nieprzywiedlnego łańcucha Markowa
istnieją stałe \(c > 0\), \(0 < \delta < 1\) takie, że dla każdego stanu \(z\) i dla każdego \(m \in \mathbb{N}_1\) mamy:
\[
	\mathbb{P}(T_z > m) \leq c \, \delta^m.
\]

\textbf{Dowód.}
Z nieprzywiedlności i nieokresowości (korzystam z tego założenia bo Krystian pozwolił) łańcucha Markowa otrzymujemy, że istnieje \(N \in \mathbb{N}\) takie, że:
\[
	\forall n \geq N, \quad \forall x, y \quad P_{x,y}(n) > 0.
\]
W szczególności:
\[
	0 < a := \min_{x,y} P_{x,y}(N).
\]

Rozważmy teraz:
\[
	\mathbb{P}(T_z > m)
	= \mathbb{P}(X_0 \neq z, X_1 \neq z, \dots, X_m \neq z)
	\leq \mathbb{P}(X_0 \neq z, X_N \neq z, X_{2N} \neq z, \dots, X_{\lfloor m/N \rfloor N} \neq z).
\]
A zatem:
\[
	\mathbb{P}(T_z > m) \leq (1 - a)^{\lfloor m/N \rfloor}
	\leq (1 - a)^{\frac{m}{N} - 1}
	= \frac{1}{1-a} \left( (1 - a)^{1/N} \right)^m.
\]

Zatem otrzymujemy dokładnie poszukiwaną postać:
\[
	\mathbb{P}(T_z > m) \leq c \, \delta^m,
\]
gdzie:
\[
	c = \frac{1}{1-a}, \qquad \delta = (1-a)^{1/N}.
\]
